\paragraph{\trCone。}\label{sub:cones}
设 $Y_0,Y_1$ 是 $\scr{A}$ 的\trObject，而 $c\in\hom^0_{\scr{A}}(Y_0,Y_1)$ 是一个%
零次\trCocycle（$\mu^1_{\scr{A}}(c)=0$）。$c$ 的\textbf{抽象\trMappingCone}
（abstract mapping cone）是 \trAinfty-模 $\scr{C}=\scr{C}one(c)$，其定义为
\begin{equation}\label{eq:cone-definition}
\begin{aligned}
&\scr{C}(X)=\hom_{\scr{A}}(X,Y_0)[1]\oplus\hom_{\scr{A}}(X,Y_1),\\
&\mu^d_{\scr{C}}\bigl((b_0,b_1),a_{d-1},\ldots,a_1\bigr)\\
&\hspace{1.5em}=\bigl(\mu^d_{\scr{A}}(b_0,a_{d-1},\ldots,a_1),\
\mu^d_{\scr{A}}(b_1,a_{d-1},\ldots,a_1)\\
&\hspace{4em}+\mu^{d+1}_{\scr{A}}(c,b_0,a_{d-1},\ldots,a_1)\bigr).
\end{aligned}
\end{equation}
直接计算表明\eqref{eq:module-associativity-equation} 得到满足。至于 $\scr{C}$
的\trCunital{}性，需要验证的是：若 $e_X\in\hom^0_{\scr{A}}(X,X)$ 代表
$H(\scr{A})$ 中的\trIdentityMorphism，则由
$(b_0,b_1)\mapsto\mu^2_{\scr{C}}((b_0,b_1),e_X)$ 诱导的
$H(\scr{C}(X),\mu^1_{\scr{C}})$ 的\trEndomorphism{}是恒等。我们已经知道它是%
\trIdempotent{}的，而考察 $\scr{C}(X)$ 中 $b_0=0$ 的子\trComplex{}及其%
\trQuotient{}即可看出，它还是\trInvertible{}的，从而就是恒等。我们把%
\trQuasiRepresent{}该抽象\trMappingCone{}的\trObject{}记作
$\mathrm{Cone}(c)$，并简称为 $c$ 的\trMappingCone。照例，这样的%
\trObject（如果存在的话）在 $H^0(\scr{A})$ 中%
\trCanonical{}\trIsomorphism{}的意义下唯一。

$\scr{C}one(c)$ 如何依赖于\trCocycle{} $c$ 在一个固定\trCohomology{}类内的
选取，考察这一点颇有启发。若取另一个代表元 $\tilde{c}=c+\mu^1_{\scr{A}}(h)$，
则
\begin{equation}\label{eq:cone-change-of-cocycle}
\begin{aligned}
&t^1(b_0,b_1)=\bigl((-1)^{|b_0|-1}b_0,\
(-1)^{|b_1|}b_1+\mu^2_{\scr{A}}(h,b_0)\bigr),\\
&t^d\bigl((b_0,b_1),a_{d-1},\ldots,a_1\bigr)
=\bigl(0,\ \mu^{d+1}_{\scr{A}}(h,b_0,a_{d-1},\ldots,a_1)\bigr)
\quad\text{其中 }d\ge 2
\end{aligned}
\end{equation}
定义出一个\trModule\trHomomorphism{} $\scr{C}one(c)\to\scr{C}one(\tilde{c})$，
由\cref{thm:module-composition-quasiisomorphism}，它是 $H^0(\scr{Q})$ 中的%
\trIsomorphism（此处以及下文的类似情形中，$|b_0|$ 都指 $b_0$ 在
$\hom_{\scr{A}}(X,Y_0)$ 中的\trDegree，也就是在施加\trSuspension{} $[1]$ 之前的%
\trDegree）。因此，\trCohomology{}类 $[c]$ 在该\trCategory{}中把 $\scr{C}one(c)$
确定到相差一个\emph{非\trCanonical}的\trIsomorphism（之所以非\trCanonical，
是因为\eqref{eq:cone-change-of-cocycle} 需要选定 $h$）。按定义，对作为
$H^0(\scr{A})$ 中\trObject{}的 $\mathrm{Cone}(c)$ 而言也有同样的结论。

\begin{remark}\label{thm:cone-shift-sum-relation}
在\trDirectSum、\trSuspension{}与\trMappingCone{}之间存在一个初等的关系：
平凡\trMorphism{} $c=0\in\hom_{\scr{A}}(Y_0,Y_1)$ 的\trCone{}\trQuasiRepresent
$SS\scr{Y}_0\oplus\scr{Y}_1$。比如说，假设 $\scr{A}$ 对\trMappingCone{}封闭，
即 $\mathrm{Cone}(c)$ 对一切 $c$ 都存在。先取一个辅助\trObject{} $Z$ 和某个
代表恒等的链 $e_Z$，便可发现 $\scr{C}one(e_Z)$ \trIsomorphic{}于零模。于是，
对任意 $Y$，零映射 $Y\to\mathrm{Cone}(e_Z)$ 的\trCone{}其实就是 $SY$，这说明
$\scr{A}$ 对\trSuspension{} $S^{\sigma}$（$\sigma>0$）封闭（显然，由此还可得知
$SY_0\oplus SY_1$ 对一切 $Y_0,Y_1$ 都存在）。
\end{remark}
